% ----------------------------------------------------------------------
% Tool lemma: constant-envelope upper semicontinuity of OCE risk measures.
% This file is an *artifact* for Section 4.3 of the thesis (tools for the
% repair). It repairs failure point (F3): the limit passage inside the
% risk measure along almost surely convergent sequences, provided a
% deterministic lower envelope is available.
% When integrating into the final thesis, strip the standalone preamble
% and include this under the same numbering as the surrounding document
% (keep the "% draft: ..." comment to preserve the original label).
% Cross-references below are faked (hardcoded) so this artifact compiles
% without defining the targets; comments say what they should point to.
% ----------------------------------------------------------------------
\documentclass[11pt,a4paper]{report}
\usepackage[T1]{fontenc}
\usepackage[utf8]{inputenc}
\usepackage{amsmath,amssymb,amsthm}
\usepackage{enumitem}

\numberwithin{equation}{chapter}

\theoremstyle{plain}
\newtheorem{proposition}{Proposition}[chapter]
\newtheorem{lemma}[proposition]{Lemma}
\newtheorem{theorem}[proposition]{Theorem}
\newtheorem{corollary}[proposition]{Corollary}
\theoremstyle{definition}
\newtheorem{assumption}[proposition]{Assumption}
\newtheorem{definition}[proposition]{Definition}
\newtheorem{example}[proposition]{Example}
% Upright body, bold head (not the italic amsthm "remark" style).
\newtheorem{remark}[proposition]{Remark}

\newcommand{\E}{\mathbb{E}}
\newcommand{\PP}{\mathbb{P}}
\newcommand{\R}{\mathbb{R}}
\newcommand{\N}{\mathbb{N}}
\newcommand{\F}{\mathcal{F}}
\newcommand{\HH}{\mathcal{H}}
\newcommand{\Hu}{\mathcal{H}^{\mathrm{u}}}
\newcommand{\NNet}{\mathcal{N\!N}}
\newcommand{\as}{\ \PP\text{-a.s.}}

\begin{document}

\setcounter{chapter}{3}

\chapter{Artifact: constant-envelope upper semicontinuity}

\subsection*{A tool for the limit passage inside the risk measure}

% FAKE REF: "(F3)" -> the failure-point label of Section 4.2 of the
%   thesis (see prop4_3_finite.tex): continuity of rho along a.s.
%   convergence fails on general Omega.
The failure point (F3) is the passage to the limit inside the risk
measure: on a general probability space, a convex risk measure need not
be upper semicontinuous along almost surely convergent sequences. The
next lemma restores exactly this passage for OCE risk measures, provided
all objectives share a deterministic lower envelope. It is the tool used
in the proof of the approximation theorem.

% draft: Lemma 4.4 (constant-envelope upper semicontinuity of OCEs);
% same statement as lem:usc-A in extension_with_A1_to_A6.tex.
\begin{lemma}\label{lem:usc}
  Let $\ell:\R\to\R$ be convex and non-decreasing and let $\rho_\ell$ be
  the associated OCE. Let $(Y_m)_{m\in\N}$ and $Y$ be random variables
  with
  \[
    Y_m\ge-B\as\text{ for all }m,
  \]
  for some deterministic $B\ge0$. If $\liminf_{m\to\infty}Y_m\ge Y$
  $\PP$-a.s., then
  \[
    \limsup_{m\to\infty}\rho_\ell(Y_m)\le\rho_\ell(Y).
  \]
\end{lemma}

\begin{proof}
  Fix $w\in\R$. Since $-(Y_m+w)\le B+|w|$ and $\ell$ is non-decreasing,
  \[
    g_w:=\ell(B+|w|)\ge\ell\bigl(-(Y_m+w)\bigr),
  \]
  so $g_w$ is a deterministic (hence integrable) constant dominating
  $\ell\bigl(-(Y_m+w)\bigr)$. Extend $\ell$ to
  $[-\infty,\infty)$ by $\ell(-\infty):=\lim_{t\to-\infty}\ell(t)$; the
  extension is continuous and non-decreasing. For any real sequence
  $(a_m)$ with $\sup_m a_m<\infty$ one has
  $\limsup_m\ell(a_m)\le\ell(\limsup_m a_m)$. Applying this with
  $a_m=-(Y_m+w)$ and using
  $\limsup_m\bigl(-(Y_m+w)\bigr)=-\liminf_m(Y_m+w)\le-(Y+w)$ a.s.\ yields
  \[
    \limsup_{m\to\infty}\ell\bigl(-(Y_m+w)\bigr)\le\ell\bigl(-(Y+w)\bigr)\as.
  \]
  Reverse Fatou applied to the sequence
  \[
    g_w\ge\ell\bigl(-(Y_m+w)\bigr)
  \]
  gives
  \[
    \limsup_{m\to\infty}\E\bigl[\ell\bigl(-(Y_m+w)\bigr)\bigr]
    \le\E\bigl[\ell\bigl(-(Y+w)\bigr)\bigr].
  \]
  Since $\rho_\ell(Y_m)\le w+\E[\ell(-(Y_m+w))]$, one obtains
  $\limsup_m\rho_\ell(Y_m)\le w+\E[\ell(-(Y+w))]$ for every $w$, and
  taking the infimum over $w$ completes the proof.
\end{proof}

\end{document}
